\section*{Hoofdstuk \ref{ch:b2:plant-flowering} --- Generatieve ontwikkeling en de regeling van de bloei}

\begin{solution}{exo:b2:plant-flowering:1}
Vegetatief \omterm{def:b2:plant-meristems:meristem}{meristeem}: maakt onbeperkt bladeren met okselknoppen
(onbegrensd). \omterm{def:b2:plant-flowering:transition}{Bloeiwijzemeristeem}: maakt schutbladen en in hun oksels
bloemmeristemen, bij veel soorten onbegrensd. Bloemmeristeem: maakt de vier
kransen en stopt --- begrensd.
\end{solution}

\begin{solution}{exo:b2:plant-flowering:2}
\omterm{prop:b2:plant-flowering:photoperiod}{Korte-dagplanten} bloeien wanneer de dag korter is dan een kritische lengte
(chrysant, sojaboon, rijst); \omterm{prop:b2:plant-flowering:photoperiod}{lange-dagplanten} wanneer hij langer is
(spinazie, tarwe, \emph{Arabidopsis}); dagneutrale planten bloeien wanneer
ze een bepaalde grootte bereiken (tomaat, maïs). Een \omterm{prop:b2:plant-flowering:photoperiod}{korte-dagplant} meet de
lengte van de ononderbroken nacht.
\end{solution}

\begin{solution}{exo:b2:plant-flowering:3}
Wildtype: kelkblad, kroonblad, \omterm{def:b2:angiosperm-reproduction:flower}{meeldraad}, \omterm{def:b2:angiosperm-reproduction:flower}{vruchtblad}. Zonder A: \omterm{def:b2:angiosperm-reproduction:flower}{vruchtblad},
\omterm{def:b2:angiosperm-reproduction:flower}{meeldraad}, \omterm{def:b2:angiosperm-reproduction:flower}{meeldraad}, \omterm{def:b2:angiosperm-reproduction:flower}{vruchtblad}. Zonder B: kelkblad, kelkblad, \omterm{def:b2:angiosperm-reproduction:flower}{vruchtblad},
\omterm{def:b2:angiosperm-reproduction:flower}{vruchtblad}. Zonder C: kelkblad, kroonblad, kroonblad, kelkblad, en de \omterm{def:b2:angiosperm-reproduction:flower}{bloem}
herhaalt zich binnenin.
\end{solution}

\begin{solution}{exo:b2:plant-flowering:4}
Weken kou en vocht (de winter is voorbij); licht (het \omterm{def:b2:angiosperm-reproduction:seed}{zaad} ligt dicht bij
het oppervlak); wisselende temperaturen (kale grond, geen bladerdak);
uitspoeling door regen (de regens zijn gekomen); opruwing (doorgang door een
darm of schuren in de bodem); rook en hitte (een brand heeft de grond
vrijgemaakt).
\end{solution}

\begin{solution}{exo:b2:plant-flowering:5}
16/8: nacht van 8 uur, onder de kritische 10: nee. 12/12: ja.
12 licht, 6 donker, flits, 6 donker: twee nachten van 6 uur: nee.
8 licht, 4 donker, 4 licht, 8 donker: langste nacht 8 uur: nee.
\end{solution}

\begin{solution}{exo:b2:plant-flowering:6}
\omterm{def:b2:plant-flowering:florigen}{Florigeen} is een via enten overdraagbaar signaal dat in het blad wordt
gemaakt en bij korte- en lange-dagsoorten hetzelfde is --- het verschil zit
in wanneer het blad het maakt, niet in wat het maakt. Met onderbroken floëem
induceert de ent geen bloei: het signaal reist via het floëem.
\end{solution}

\begin{solution}{exo:b2:plant-flowering:7}
$n = \ln(1/0.15)/0.05 = 1.90/0.05 = 38$ koude dagen. De periodes zijn samen
$47$ dagen: de plant is gevernaliseerd op de 11e dag van de derde periode.
Met $\alpha = 0.02$ zou ze 95 dagen nodig hebben en zou ze het niet zijn.
\end{solution}

\begin{solution}{exo:b2:plant-flowering:8}
Het gen van de repressor wordt stilgelegd door chromatinemarkeringen die bij
elke mitose worden gekopieerd, zodat elke cel van de scheut de winter
onthoudt; in de kiembaan en het vroege embryo worden de markeringen gewist en
komt het gen weer tot expressie. Nuttig, omdat een \omterm{def:b2:angiosperm-reproduction:seed}{zaad} dat in de zomer
wordt uitgestrooid, de toestemming om te bloeien niet van zijn ouder mag
erven: het moet op een eigen winter wachten.
\end{solution}

\begin{solution}{exo:b2:plant-flowering:9}
R: kiemen (fytochroom omgezet in de verroodabsorberende vorm, Pfr, de
actieve). R dan VR: niet (teruggezet naar Pr). R, VR, R: kiemen. R, VR, R,
VR: niet. Alleen de laatste flits telt, omdat elke omzetting volledig is en
het \omterm{def:b2:angiosperm-reproduction:seed}{zaad} de eindtoestand van het pigment afleest.
\end{solution}

\begin{solution}{exo:b2:plant-flowering:10}
$p = 0.7$: $r = 1.28$, $1.42$, $1.40$, $-\infty$ voor $g = 0.4, 0.6, 0.8,
1$: het beste bij ongeveer $g = 0.7$. $p = 0.95$: $1.99$, $2.33$, $2.55$,
$-\infty$: het optimum schuift naar $g = 1$ (ongeveer $0.95$), met nog maar
een symbolische reserve. Woestijnen, waar goede jaren zeldzaam zijn,
begunstigen een zware \omterm{def:b2:angiosperm-reproduction:seed}{kiemrust} en een grote zaadbank; weilanden, waar de
meeste jaren goed zijn, begunstigen het laten kiemen van bijna alles.
\end{solution}

\begin{solution}{exo:b2:plant-flowering:11}
B in alle kransen: kroonblad, kroonblad, \omterm{def:b2:angiosperm-reproduction:flower}{meeldraad}, \omterm{def:b2:angiosperm-reproduction:flower}{meeldraad}. C in alle
kransen (C onderdrukt A): \omterm{def:b2:angiosperm-reproduction:flower}{vruchtblad}, \omterm{def:b2:angiosperm-reproduction:flower}{meeldraad}, \omterm{def:b2:angiosperm-reproduction:flower}{meeldraad}, \omterm{def:b2:angiosperm-reproduction:flower}{vruchtblad}.
Zonder A en zonder B: overal alleen C --- \omterm{def:b2:angiosperm-reproduction:flower}{vruchtbladen} in alle vier de
kransen. C beëindigt ook het \omterm{def:b2:plant-meristems:meristem}{meristeem}; zonder C blijft het centrum
onbegrensd en maakt het \omterm{def:b2:angiosperm-reproduction:flower}{bloem} in \omterm{def:b2:angiosperm-reproduction:flower}{bloem}.
\end{solution}

\begin{solution}{exo:b2:plant-flowering:12}
De daglengte wordt over de nachten geïntegreerd door het samenvallen van
licht met de klok, en de inductie vraagt meerdere zulke nachten; kou wordt
geïntegreerd als opstapelende chromatinemarkeringen; de zaadbank integreert
over de jaren door de kieming te spreiden. Een integrator negeert één
afwijkende dag, een dooi in januari of één slecht jaar; een drempel op de
ogenblikkelijke waarde zou de plant in een warme week in februari laten
bloeien, of elk \omterm{def:b2:angiosperm-reproduction:seed}{zaad} bij de eerste regen laten kiemen.
\end{solution}

\begin{solution}{pb:b2:plant-flowering:1}
\textbf{1.} Een nacht van 8 uur is korter dan de kritische 9.5 uur: geen
bloei. Om in augustus te verkopen, moet de kweker de nachten kunstmatig
verlengen met zwart doek.
\textbf{2.} Een nacht van 12 uur. Bloei op 15 augustus betekent dat de
inductie 8 weken eerder voltooid was, rond 20 juni; de twaalf inductieve
nachten moeten rond 8 juni beginnen.
\textbf{3.} Belicht het gewas midden in elke nacht enkele minuten: een nacht
van 14.5 uur wordt er twee van ongeveer 7 uur, beide onder de kritische
lengte. Een flits volstaat omdat fytochroom binnen minuten wordt omgezet en
de plant alleen de lengte van de ononderbroken duisternis meet.
\textbf{4.} Een flits om 21:00 laat een donkere periode over van
21:00 tot 07:30, 10.5 uur --- langer dan 9.5: dat deel induceert. De
onderbreking moet zo vallen dat geen enkel deel de kritische lengte
overschrijdt.
\textbf{5.} Eén gemiste nacht verbreekt de reeks van twaalf opeenvolgende
inductieve nachten: de telling begint opnieuw en de bloei wordt ongeveer
evenveel dagen uitgesteld als er waren opgebouwd. Een stekelnoot heeft één
enkele inductieve nacht nodig en zou al vastgelegd zijn geweest.
\textbf{6.} Fytochroom absorbeert geen groen licht, dus de plant ziet de
flits niet; rood zet het om in de actieve vorm, die als dag wordt gelezen;
verrood direct daarna zet het terug, zodat de nacht niet wordt onderbroken.
\textbf{7.} De \omterm{prop:b2:plant-flowering:photoperiod}{lange-dagplant} bloeit: \omterm{def:b2:plant-flowering:florigen}{florigeen} steekt de entverbinding over
via het floëem. Het identificeert een beweeglijk, via enten overdraagbaar
signaal dat beide daglengteklassen gemeen hebben.
\textbf{8.} $n = \ln 10/0.06 = 38.4$: 39 koude dagen.
\textbf{9.} 10 in november, 35 aan het eind van januari, en de 39e koude dag
valt op de vierde koude dag van februari: gevernaliseerd begin februari.
\textbf{10.} Zes dagen: $f = e^{-0.36} = 0.70 > 0.10$: ze blijft de hele
zomer vegetatief en geeft geen graan --- wintertarwe moet in de herfst
worden gezaaid.
\textbf{11.} Zonder de repressor bloeit hij zonder kou: hij kan in het
voorjaar worden gezaaid en dezelfde zomer worden geoogst --- een zomertarwe.
\textbf{12.} De stilleggende markeringen op het chromatine van de repressor
worden bij elke mitose gekopieerd, zodat de top het voorjaar door onthoudt;
bij de \omterm{def:b2:meiosis-heredity:meiosis}{meiose} en in het embryo worden de markeringen gewist, en moet elke
generatie haar eigen winter tellen.
\textbf{13.} Het geheugen moet zitten in de cellen die de \omterm{def:b2:angiosperm-reproduction:flower}{bloem} zullen maken
en die de winter doorstaan --- het \omterm{def:b2:plant-meristems:meristem}{meristeem} --- en het is een toestand van
hun chromatine, geen beweeglijk signaal; de bladeren die de daglengte meten,
worden afgeworpen, en voeren in plaats daarvan een eiwit uit.
\textbf{14.} Klasse C: kelkblad, kroonblad, kroonblad, kelkblad (en de \omterm{def:b2:angiosperm-reproduction:flower}{bloem}
herhaalt zich binnenin, wat de extra kroonbladen in het midden geeft).
\textbf{15.} Geen meeldraden, dus geen \omterm{def:b2:plant-life-cycles:heterospory}{stuifmeel}; waar C zwak actief is,
vormen zich enkele \omterm{def:b2:angiosperm-reproduction:flower}{vruchtbladen}, zodat er af en toe \omterm{def:b2:angiosperm-reproduction:seed}{zaad} wordt gezet.
\textbf{16.} Klasse A: \omterm{def:b2:angiosperm-reproduction:flower}{vruchtblad}, \omterm{def:b2:angiosperm-reproduction:flower}{meeldraad}, \omterm{def:b2:angiosperm-reproduction:flower}{meeldraad}, \omterm{def:b2:angiosperm-reproduction:flower}{vruchtblad}.
\textbf{17.} C onderdrukt A: krans 1 wordt \omterm{def:b2:angiosperm-reproduction:flower}{vruchtblad}, krans 2 (B en C)
\omterm{def:b2:angiosperm-reproduction:flower}{meeldraad} --- \omterm{def:b2:angiosperm-reproduction:flower}{vruchtblad}, \omterm{def:b2:angiosperm-reproduction:flower}{meeldraad}, \omterm{def:b2:angiosperm-reproduction:flower}{meeldraad}, \omterm{def:b2:angiosperm-reproduction:flower}{vruchtblad}, de \omterm{def:b2:angiosperm-reproduction:flower}{bloem} van de
A-mutant.
\textbf{18.} De ABC-genen vormen één familie van \omterm{prop:b2:cell-differentiation:transcriptional}{transcriptiefactoren},
geërfd van de gemeenschappelijke voorouder van alle bloemplanten; de module
bestond al voordat rozen, leeuwenbekken en \emph{Arabidopsis} uiteengingen,
en elk heeft haar behouden.
\textbf{19.} B zonder A of C specificeert geen bloemorgaan: een bladachtig
of schutbladachtig orgaan.
\textbf{20.} $p = 0.8$: $r = 1.44$ ($g = 0.5$), $1.59$ ($0.7$), $1.55$
($0.9$).
\textbf{21.} Van de drie $g = 0.7$ (het echte optimum ligt bij ongeveer
$0.8$); bij $g = 1$ is de factor in een slecht jaar nul en wordt de lijn
door het eerste slechte jaar uitgedoofd.
\textbf{22.} $p = 0.5$: $0.51$, $0.41$, $-0.03$: $g = 0.5$ is de beste,
en een nog lagere zou beter zijn --- meer \omterm{def:b2:angiosperm-reproduction:seed}{kiemrust} in de woestijn.
\textbf{23.} De zaden van de teler worden opgeslagen, niet begraven, en
overleven met zekerheid; hij hoeft binnen het veld geen risico te spreiden,
maar moet in de schuur een reserve houden om na een mislukt jaar opnieuw te
zaaien --- dezelfde reserve, maar door hem aangehouden in plaats van door de
bodem.
\textbf{24.} Een klein \omterm{def:b2:angiosperm-reproduction:seed}{zaad} heeft reserves voor een scheut van een of twee
centimeter; kiemend in het donker onder de grond zou het sterven voordat het
het licht bereikte, dus wacht het op licht --- dat het alleen dicht bij het
oppervlak bereikt. Ploegen brengt begraven zaden naar boven en veroorzaakt
een golf onkruid.
\textbf{25.} Kritische nacht 9.5 uur, bedekken vanaf ongeveer 8 juni voor
een bloei op 15 augustus; 39 koude dagen; beste $g \approx 0.8$ voor
$p = 0.8$ en ongeveer $0.5$ voor $p = 0.5$.
\end{solution}
